%==============================================================================
% Section 4: Mode Dictionary: T^3 x S^1 under EC-Weyl
%==============================================================================
\section{Mode Dictionary: $\Tthree\times\Sone$ under EC--Weyl}
\label{sec:t3}

This section analyses the effect of the EC--Weyl coupling on the homogeneous
geometric degrees of freedom of $\Tthree\times\Sone$.  The flatness of $\Tthree$
($C^a_{bc}=0$) causes a triple triviality, and we establish Theorem~5
(torsion--volume coupling) as the origin of the squash--shear cross-term.

\subsection{$\Tthree$ Flatness Theorem}
\label{sec:t3:flat}

$\Tthree = \mathbb{R}^3/\mathbb{Z}^3$ is a flat torus with
\begin{equation}
  C^a_{bc}(\Tthree) = 0 \quad \forall\,a,b,c.
\end{equation}
This flatness causes the Weyl tensor to vanish, setting all spin-2/spin-1
masses to zero.

\paragraph{Spin-2 mass (squash direction).}
By exact SymPy computation,
\begin{equation}
  \frac{\partial^2 V}{\partial\varepsilon^2}\bigg|_{\varepsilon=0}
  = 0 \quad \forall\,\eta, V, \alpha.
\end{equation}
With $C^a_{bc}=0$, the Weyl tensor vanishes and no geometric mass term for spin-2
arises.  Script: \texttt{paper03ec/t3\_mode\_dictionary.py}.

\paragraph{Spin-2 mass (shear direction).}
Similarly,
\begin{equation}
  \frac{\partial^2 V}{\partial s^2}\bigg|_{s=0}
  = 0 \quad \forall\,\eta, V, \alpha
\end{equation}
(SymPy exact zero; script: \texttt{proofs/t3\_flatness\_null\_test.py}).

\paragraph{Spin-1 mass.}
With $C^a_{bc}=0$, all spin-1 components are massless: $m^2_\omega = 0$.

\subsection{Mode Dictionary Table ($\Tthree\times\Sone$)}
\label{sec:t3:table}

\begin{table}[H]
\centering
\small
\begin{tabular}{lccc}
\toprule
Quantity & AX ($V=0,\,\eta\neq 0$) & VT ($\eta=0,\,V\neq 0$) & MX ($V\neq 0,\,\eta\neq 0$) \\
\midrule
$C^2_\EC$ & $=C^2_\LC=0$ & $=C^2_\LC=0$ & $16V^2\eta^2/(3r^2)$ \\
AX/VT dropout & \checkmark & \checkmark & $\times$ \\
spin-0: $r$ & propagating & propagating & propagating \\
spin-0: $\eta,V$ & $\eta$ non-prop., $V=0$ & $\eta=0$, $V$ non-prop. & both non-prop. \\
spin-2 mass $\mu^2$ & $0$ ($\Tthree$ flatness) & $0$ ($\Tthree$ flatness) & $0$ ($\Tthree$ flatness) \\
spin-1 mass $m^2_\omega$ & $0$ (massless) & $0$ & $0$ \\
\bottomrule
\end{tabular}
\caption{$\Tthree\times\Sone$ mode dictionary under EC--Weyl.}
\label{tab:t3}
\end{table}

\subsection{Triple Triviality of $\Tthree$}
\label{sec:t3:triple}

The effect of the EC--Weyl coupling in $\Tthree\times\Sone$ is trivial in three
senses.

\paragraph{(i) AX dropout.}
From $C^a_{bc}=0$ we get $C^2_\LC=0$; Theorem~\ref{thm:dropout} (AX dropout)
then gives $C^2_\EC = C^2_\LC = 0$.  The coupling $\alpha$ does not modify the
effective potential in the AX background.

\paragraph{(ii) VT dropout.}
Similarly, $\eta=0$ gives Theorem~\ref{thm:dropout} (VT dropout) and
$C^2_\EC=0$.  In the VT background too, $\alpha$ is inactive.

\paragraph{(iii) MX correction.}
In the MX background ($V\eta\neq 0$) where AX/VT dropout does not apply,
$C^2_\EC = 16V^2\eta^2/(3r^2)\neq 0$.  For $\alpha<0$ this correction gives
$\Delta\Veff^\EC > 0$, pushing the potential upward.  Whether a finite MX
critical point exists is examined algebraically in \S\ref{sec:t3:mx}.

The auxiliary EFT coefficient is also
\begin{equation}
  C_\delta(\Tthree) = 0,
  \qquad
  \frac{\partial C_\delta(\Tthree)}{\partial\alpha} = 0
\end{equation}
(Theorem~\ref{thm:delta}, Appendix~\ref{app:D}).  Hence in $\Tthree$, neither
the physical sector nor the auxiliary sector develops a new homogeneous cubic
channel from the EC extension.

\subsection{Theorem~5: Torsion--Volume Coupling}
\label{sec:t3:thm5}

\begin{theorem}[Torsion--volume coupling]\label{thm:tvol}
  In $\Tthree\times\Sone$, the mixed second derivative of the squash and shear
  parameters satisfies
  \begin{equation}
    \frac{\partial^2\Veff}{\partial\varepsilon\,\partial s}\bigg|_{\varepsilon=s=0}
    = \frac{48\pi^4 L\eta^2 r}{\kappa^2} \neq 0 \quad (\eta\neq 0).
  \end{equation}
  The origin of this cross-term is the kinematic mixing of the volume element
  $(1+\varepsilon)(1+s)$, independent of the Weyl mass.
\end{theorem}

\begin{proof}[Proof outline]
  The effective potential of $\Tthree$ takes the form (AX background, $V=0$):
  \begin{equation}
    V_\mathrm{eff}^{T^3}(\varepsilon,s)
    \supset \frac{16\pi^4 Lr}{\kappa^2}\cdot 9\eta^2\cdot(1+\varepsilon)(1+s)
    + \cdots
  \end{equation}
  (with $C^a_{bc}=0$ the Weyl contribution is absent, and the torsion kinetic
  term $\propto\eta^2$ is multiplied by the volume factor $(1+\varepsilon)(1+s)$).
  The mixed derivative yields:
  \begin{equation}
    \frac{\partial^2 V_\mathrm{eff}^{T^3}}{\partial\varepsilon\,\partial s}
    \bigg|_{\varepsilon=s=0}
    = \frac{16\pi^4 Lr}{\kappa^2}\cdot 9\eta^2\cdot 1
    = \frac{48\pi^4 L\eta^2 r}{\kappa^2}.
  \end{equation}

  \textit{Confirmation that this is not a Weyl mass.}  In the AX background
  ($V=0$), $C^a_{bc}=0$ gives $C^2_\LC=0$, and Theorem~\ref{thm:dropout} (AX
  dropout) gives $C^2_\EC = 0$, so this cross-term is $\alpha$-independent and
  carries no Weyl-mass contribution.  Numerical check (AX background): the
  cross-term value is identical for $\alpha=0$ and $\alpha=-1$ (SymPy exact).
  In the MX background ($V\neq 0$), $\alpha$-dependent terms join the
  cross-term via $C^2_\EC=16V^2\eta^2/(3r^2)\neq 0$, but the spin-2 null test
  ($\partial^2 V/\partial\varepsilon^2=0$, $\partial^2 V/\partial s^2=0$)
  continues to hold (because $C^2_\EC$ is independent of $\varepsilon,s$).
  Scripts: \texttt{proofs/t3\_flatness\_null\_test.py},
  \texttt{paper03ec/squash\_shear\_cross\_term.py}.
\end{proof}

\subsection{Absence of MX Critical Points}
\label{sec:t3:mx}

We examine critical points of the $\Tthree$ MX sector for $\alpha<0$.
Solving $\partial\Veff/\partial\eta = \partial\Veff/\partial V = 0$ symbolically
by SymPy, the MX solution is
\begin{equation}
  V^* = \pm\frac{3}{4}\sqrt{\frac{1}{\alpha}},
  \qquad
  \eta^* = \pm\frac{r}{4}\sqrt{\frac{1}{\alpha}},
\end{equation}
which is purely imaginary for $\alpha<0$ (no real solution).  Therefore
\textbf{no finite real MX critical point exists in $\Tthree$ for $\alpha<0$}.
This is consistent with the analysis in \S\ref{sec:t3:triple} that the MX
correction pushes the potential upward, giving no new homogeneous local minimum
in the MX sector of $\Tthree$.  Script: \texttt{paper03ec/t3\_mx\_vacuum\_search.py}.
